\documentclass[acmsmall,screen,review,anonymous]{acmart}

\citestyle{acmauthoryear}

%% Packages
\usepackage{amsmath,amssymb}
\usepackage{mathpartir}

%% Agda-style code
\usepackage{listings}
\lstset{
  basicstyle=\ttfamily\small,
  keywordstyle=\bfseries,
  columns=flexible,
  mathescape=true,
  literate={→}{$\to$}1 {∀}{$\forall$}1 {λ}{$\lambda$}1
           {≡}{$\equiv$}1 {×}{$\times$}1 {Σ}{$\Sigma$}1
           {⟧}{$\rrbracket$}1 {⟦}{$\llbracket$}1
}

\begin{document}

\title{Authenticated Data Structures, for Free}
\subtitle{An ICFP Pearl}

\author{Anonymous}
\affiliation{\institution{Anonymous}}

\begin{abstract}
Miller et al.\ (2014) built authenticated data structures---Merkle
trees, authenticated dictionaries, and the like---as a generic
library, with prover and verifier as two instantiations of an abstract
\emph{auth kit} interface.  We formalize the security of this approach
in Agda, with two results.

The main theorem is a \emph{collision extraction}: if an adversary
provides a proof stream that makes the verifier accept an incorrect
result, we constructively produce a collision in the underlying hash
function.  This is a standard cryptographic reduction, but the proof
is remarkably simple---induction on the computation tree (a free monad
of hash-checked lookups), 129 lines of plain Agda, with no
assumptions beyond collision resistance.

As a supporting result, we formalize Atkey's (2016) claim that
parametric polymorphism over the auth kit guarantees correctness of
honest computations, using \emph{internal parametricity}
(agda-bridges) in Cubical Agda.
\end{abstract}

\maketitle

% TODO: Introduction
% - Miller et al. 2014 built ADS as a library
% - Intuition: polymorphism = security, but couldn't close the gap
% - Two contributions: (1) collision extraction — the real security theorem,
%   (2) Atkey 2016 parametricity claim formalized — correctness of honest code
% - Collision extraction is simple and self-contained (plain Agda)
% - Parametricity formalization uses agda-bridges (internal parametricity)

% TODO: Background
% - ADS basics (Merkle trees, prover/verifier)
% - Collision resistance (not injectivity!)
% - Brief on parametricity / agda-bridges (for Section 5)

% ============================================================
\section{The Auth Kit}\label{sec:auth-kit}
% ============================================================

An authenticated data structure scheme consists of two parties---a
\emph{prover} who maintains the data and produces Merkle-style
certificates, and a \emph{verifier} who checks them.  The key
observation of Miller et al.~\cite{miller2014} is that both can be
obtained from a single generic program, written once against an
abstract interface we call the \emph{auth kit}.

The kit has two layers.  The first is a monadic computation
environment: a type constructor $\mathsf{m}$ with $\mathsf{ret}$ and
$\mathsf{bind}$.  Any program can be written in this environment---it
is just a monad.  The second layer adds two operations, $\mathsf{auth}$
and $\mathsf{unauth}$, that mark \emph{authenticated boundaries}: places
where data crosses between a trusted local representation and an
untrusted authenticated one.

A program that never calls $\mathsf{auth}$ or $\mathsf{unauth}$ is
perfectly well-typed.  It simply does not use the authentication
machinery, and gains nothing from it---much like a program in
\texttt{IO} that never performs I/O.  The \emph{performance} benefit of
authenticated data structures comes from placing $\mathsf{auth}$ and
$\mathsf{unauth}$ judiciously: the prover hashes at $\mathsf{auth}$
sites and skips reconstruction at $\mathsf{unauth}$ sites.  Choosing
where to place these calls requires a white-box understanding of the
intended data structure---the kit interface does not dictate this.

What the kit \emph{does} guarantee, generically, is \emph{security}:
no matter how the programmer uses $\mathsf{auth}$ and $\mathsf{unauth}$,
the prover cannot cause the verifier to accept an incorrect answer.
This is the free theorem we formalize.

\subsection{The Interface}

Concretely, an auth kit bundles a monadic computation environment
together with the auth/unauth capabilities:
%
\begin{align*}
  &\textbf{Computation environment:} \\
  \mathsf{m}       &: \mathsf{Type} \to \mathsf{Type}
    && \text{(effect type)} \\
  \mathsf{ret}     &: \forall X.\; X \to \mathsf{m}\,X \\
  \mathsf{bind}    &: \forall X\,Y.\; \mathsf{m}\,X \to
                      (X \to \mathsf{m}\,Y) \to \mathsf{m}\,Y \\[4pt]
  &\textbf{Authentication capabilities:} \\
  \mathsf{au}      &: \mathsf{Type} \to \mathsf{Type}
    && \text{(authenticated wrapper)} \\
  \mathsf{auth}    &: \forall X.\; X \to \mathsf{au}\,X \\
  \mathsf{unauth}  &: \forall X.\; \mathsf{au}\,X \to \mathsf{m}\,X
\end{align*}
%
The monad $\mathsf{m}$ captures the computational side-effects
particular to each role: the prover threads Merkle-tree state, the
verifier propagates hash-checking failures, and the identity kit (below)
does nothing.  The wrapper $\mathsf{au}$ marks data as
``authenticated''---for the prover this means a value together with its
hash, for the verifier just the hash, and for the identity kit the bare
value.

The kit \emph{designer} makes two choices: a monad suited to the
role (state, error, identity, etc.), and an implementation of the
$\mathsf{auth}$/$\mathsf{unauth}$ instructions in that monad.  The
\emph{program author} writes against the abstract interface, choosing
where to place authenticated boundaries for the data structure at hand.

In Agda:
%
\begin{lstlisting}
AuthKit : Type (lsuc l)
AuthKit = $\Sigma$ (Type l $\to$ Type l) $\lambda$ m $\to$
  $\Sigma$ (Type l $\to$ Type l) $\lambda$ au $\to$
      (($\forall$ X $\to$ X $\to$ m X)                    -- ret
    $\times$ (($\forall$ X Y $\to$ m X $\to$ (X $\to$ m Y) $\to$ m Y)  -- bind
    $\times$ (($\forall$ X $\to$ X $\to$ au X)                  -- auth
    $\times$ (($\forall$ X $\to$ au X $\to$ m X))))              -- unauth
\end{lstlisting}

\subsection{The Identity Kit}

The simplest kit sets both $\mathsf{m}$ and $\mathsf{au}$ to the
identity---no effects, no wrapping:
%
\begin{lstlisting}
IdKit : AuthKit
IdKit = ($\lambda$ X $\to$ X) , ($\lambda$ X $\to$ X) ,
  ($\lambda$ X x $\to$ x) , ($\lambda$ X Y x k $\to$ k x) ,
  ($\lambda$ X x $\to$ x) , ($\lambda$ X x $\to$ x)
\end{lstlisting}
%
Running a polymorphic program with \texttt{IdKit} strips away all
authentication machinery and yields the ``pure'' result.  This is the
reference computation: the answer the verifier should get.

\subsection{Lawful Kits}

We require two laws:
%
\begin{enumerate}
\item \textbf{Monad left identity.}\;
  $\mathsf{bind}\;(\mathsf{ret}\;x)\;k \;=\; k\;x$
\item \textbf{Auth roundtrip.}\;
  $\mathsf{unauth}\;(\mathsf{auth}\;x) \;=\; \mathsf{ret}\;x$
\end{enumerate}
%
The first is a standard monad law.  The second says that wrapping a
value and immediately unwrapping it is a no-op (up to the monadic
return)---an authenticated value round-trips faithfully.  These are
precisely the laws the formalization needs: monad left identity is used
once (for $\mathsf{bind}$ compatibility) and the auth roundtrip is used
once (for $\mathsf{unauth}$ compatibility).  No further monad laws or
naturality conditions appear.

\subsection{Polymorphic Programs}

A program written against the auth kit has the type
\[
  f : \forall\, \mathit{kit}.\; A \to \mathsf{m}_{\mathit{kit}}\,A
\]
for some fixed type $A$.  The program receives a kit and may use any of
its operations.  Crucially, it is parametrically polymorphic over the
kit: it cannot inspect which implementation it received, nor behave
differently for the prover than for the verifier.  This is the
structural constraint that gives us the security theorem for free.

% ============================================================
\section{Collision Extraction}\label{sec:extraction}
% ============================================================

The security theorem for authenticated data structures is a
\emph{reduction}: if an adversary fools the verifier into accepting a
wrong answer, we extract a collision in the hash function.  We do not
assume the hash is injective---only that finding collisions is hard.

\subsection{Computation Trees}

We model a verified computation as a free monad of hash-checked
lookups:
%
\begin{lstlisting}
data Comp (R : Set) : Set where
  ret  : R $\to$ Comp R
  step : Digest $\to$ (Val $\to$ Comp R) $\to$ Comp R
\end{lstlisting}
%
A $\mathsf{ret}\;r$ returns the result $r$.  A
$\mathsf{step}\;d\;k$ says ``give me a value $v$ with
$\mathsf{hash}\;v = d$, then continue with $k\;v$.''  This is what
each $\mathsf{unauth}$ call compiles to: the digest $d$ was computed
by an earlier $\mathsf{auth}$, and the continuation $k$ captures the
rest of the program, including any future digests that may depend on
$v$.

\subsection{The Verifier}

The verifier runs a computation against a \emph{proof stream}---a list
of values provided by a prover (honest or adversarial):
%
\begin{lstlisting}
run : Comp R $\to$ List Val $\to$ Outcome R
run (ret r)    s       = ok r s
run (step d k) []      = fail
run (step d k) (v :: s) with d $\stackrel{?}{=}$ hash v
... | yes _ = run (k v) s
... | no  _ = fail
\end{lstlisting}
%
At each step, the verifier pops a value from the stream, checks that
its hash matches the expected digest, and continues.

\subsection{The Theorem}

\begin{theorem}[Collision extraction]\label{thm:extraction}
For any computation $c$, if two proof streams $s_1, s_2$ both pass
verification but produce different results, there exists a collision
in $\mathsf{hash}$:
\[
  \mathsf{run}\;c\;s_1 = \mathsf{ok}\;r_1
  \;\;\wedge\;\;
  \mathsf{run}\;c\;s_2 = \mathsf{ok}\;r_2
  \;\;\wedge\;\;
  r_1 \neq r_2
  \;\;\Longrightarrow\;\;
  \exists\, x\,y.\;\; x \neq y \;\wedge\; \mathsf{hash}\;x = \mathsf{hash}\;y
\]
\end{theorem}

\begin{proof}
By induction on the computation tree $c$.

\textbf{Case} $\mathsf{ret}\;r$: both runs return $r$, so
$r_1 = r = r_2$, contradicting $r_1 \neq r_2$.

\textbf{Case} $\mathsf{step}\;d\;k$: both streams provide values
$v_1, v_2$ that pass the hash check, so
$\mathsf{hash}\;v_1 = d = \mathsf{hash}\;v_2$.
\begin{itemize}
\item If $v_1 \neq v_2$: collision found.
\item If $v_1 = v_2$: both runs continue at the same tree node
  $k\;v_1$ with the remaining streams.  The results still differ,
  so the inductive hypothesis yields a collision.
\end{itemize}
The collision lives at the \emph{first divergence point} between
the two streams.  Before that point, both runs follow identical
paths through the tree, so the expected digest is the same---this
is what makes the hash equality valid.
\end{proof}

\subsection{Connecting to the Security Game}

The standard ADS security game is:
\begin{enumerate}
\item Fix a hash function $h$ (chosen at random from a family).
\item The honest prover runs $f$ with the prover kit, producing
  result $r$ and proof stream $\pi$.
\item An adversary (arbitrary, not parametric over the kit) produces
  a different proof stream $\pi'$.
\item The verifier runs $f$ with $\pi'$ and accepts result $r'$.
\item If $r' \neq r$, Theorem~\ref{thm:extraction} extracts a
  collision in $h$.
\end{enumerate}
The adversary need not be a program written against the auth kit
interface.  It simply provides a proof stream---a list of values.
The collision is extracted from the two streams (honest and
adversarial) together.

Note that correctness of the honest computation---that the prover
produces the right answer---is a separate property.  We address it
in Section~\ref{sec:free-theorem} via parametricity.

% ============================================================
\section{Correctness via Parametricity}\label{sec:free-theorem}
% ============================================================

\subsection{Statement}

\begin{theorem}[Purity]\label{thm:purity}
Let $\mathit{kit}$ be a lawful auth kit and
$f : \forall\,\mathit{kit}.\;A \to \mathsf{m}_{\mathit{kit}}\,A$
be parametrically polymorphic.  Then for every $x : A$:
\[
  f\;\mathit{kit}\;x \;=\;
    \mathsf{ret}_{\mathit{kit}}\!\bigl(f\;\mathit{IdKit}\;x\bigr)
\]
\end{theorem}

\begin{corollary}[Agreement]\label{cor:agreement}
For any two lawful kits $\mathit{kit}_1, \mathit{kit}_2$:
\[
  \mathit{pure}(f\;\mathit{kit}_1\;x) \;=\;
  \mathit{pure}(f\;\mathit{kit}_2\;x) \;=\;
  f\;\mathit{IdKit}\;x
\]
\end{corollary}

\begin{proof}
Apply Theorem~\ref{thm:purity} to each kit.
\end{proof}

\subsection{The Purity Relation}\label{sec:purity-rel}

The proof constructs a logical relation between an arbitrary lawful
kit and \texttt{IdKit}.  For each type constructor:

\paragraph{For $\mathsf{m}$.}
$\mathit{mx}_0 : \mathsf{m}_{\mathit{kit}}\,X_0$ relates to
$x_1 : X_1$ when there exists $x_0$ such that
$\mathsf{ret}\;x_0 = \mathit{mx}_0$ and $R\;x_0\;x_1$.
That is, $\mathit{mx}_0$ is pure.

\paragraph{For $\mathsf{au}$.}
$\mathit{ax}_0 : \mathsf{au}_{\mathit{kit}}\,X_0$ relates to
$x_1 : X_1$ when there exists $x_0$ such that
$\mathsf{auth}\;x_0 = \mathit{ax}_0$ and $R\;x_0\;x_1$.

\subsection{Compatibility}

Each operation must respect the relation.

\paragraph{$\mathsf{ret}$.}
Immediate: the witness is $x_0$ itself.

\paragraph{$\mathsf{auth}$.}
Immediate: same argument.

\paragraph{$\mathsf{bind}$.}
Given $\mathit{mx}_0 = \mathsf{ret}\;x_0$ (from the $\mathsf{m}$
relation) and a continuation $k_0$ that preserves the relation:
\[
  \mathsf{bind}\;(\mathsf{ret}\;x_0)\;k_0
  \;=\; k_0\;x_0
  \qquad\text{(monad left identity)}
\]
Then $k_0\;x_0 = \mathsf{ret}\;y_0$ (from the hypothesis on $k_0$),
so $\mathsf{bind}\;\mathit{mx}_0\;k_0 = \mathsf{ret}\;y_0$.
\emph{Uses monad left identity once.}

\paragraph{$\mathsf{unauth}$.}
Given $\mathit{ax}_0 = \mathsf{auth}\;x_0$ (from the $\mathsf{au}$
relation):
\[
  \mathsf{unauth}\;(\mathsf{auth}\;x_0) = \mathsf{ret}\;x_0
  \qquad\text{(auth roundtrip)}
\]
\emph{Uses the auth roundtrip law once.}

\subsection{Assembling the Proof}

With all four compatibility proofs, the parametricity principle
(instantiated at the purity relation) yields:
\[
  R_{\mathsf{m}}\bigl(f\;\mathit{kit}\;x,\;\;f\;\mathit{IdKit}\;x\bigr)
\]
Unfolding: $f\;\mathit{kit}\;x = \mathsf{ret}\;a$ where
$a = f\;\mathit{IdKit}\;x$.  \qed

\medskip
The two laws are precisely calibrated: each is used exactly once, in
the one compatibility proof that needs it.

% ============================================================
\section{Formalization of Correctness}\label{sec:formalization}
% ============================================================

% TODO: The NRGraph encoding in agda-bridges
% - AuthKitNRG construction (Unit,m → Unit,m,au → full kit)
% - The edge-kit-Id term (the purity relation, formalized)
% - The param call and what it returns
% - Theorem1-canonical and values-agree
% This is supporting material — the collision extraction is the main result

% ============================================================
\section{Discussion}\label{sec:discussion}
% ============================================================

% TODO:
% - The collision extraction is remarkably simple — why?
%   The free monad cleanly separates program structure from proof stream
% - Collision resistance vs injectivity: we never assume hash injectivity
% - The first-divergence-point argument: before divergence, everything identical
% - Limitations: fixed computation tree, no probabilistic reasoning,
%   collision resistance assumed not formalized as computational hardness
% - The kit pattern as a general methodology

% ============================================================
\section{Related Work}\label{sec:related}
% ============================================================

% TODO

\bibliographystyle{ACM-Reference-Format}
\bibliography{refs}

\end{document}
